\documentclass[12pt]{article}
\usepackage{amsmath}
\usepackage{amssymb}
\usepackage[margin=2.5cm]{geometry}
\usepackage{physics}
\usepackage[hidelinks]{hyperref}
\usepackage[dvipsnames]{xcolor}
\usepackage{enumitem}

\setlist[itemize]{label=--, itemsep=3pt, topsep=0.5em}

\newcommand{\uu}{\mathbf{u}}
\newcommand{\vv}{\mathbf{v}}
\newcommand{\xx}{\mathbf{x}}
\newcommand{\stress}{\boldsymbol{\sigma}}
\newcommand{\strain}{\boldsymbol{\epsilon}}
\newcommand{\shapeu}{\boldsymbol{\varphi}}
\newcommand{\shapev}{\boldsymbol{\psi}}
\newcommand{\shapef}{\boldsymbol{\tau}}
\newcommand{\llangle}{\langle\!\langle}
\newcommand{\rrangle}{\rangle\!\rangle}

\title{Élasticité linéaire 2D}
\author{Vincent Degrooff}
\date{\today}

\begin{document}
\maketitle

\tableofcontents
\newpage

\section{Modèle constitutif}

On s'intéresse à la mécanique du solide dans l'hypothèse de petites 
déformations. Les particules matérielles d'un corps soumis à différentes forces 
de volume et de surface se déplacent d'une position $\mathbf{X}$ à une position 
$\mathbf{x}(\mathbf{X},t)$. Le déplacement se note $\uu=\mathbf{x}-\mathbf{X}$ 
et satisfait les equations suivantes :

\begin{align}
    \rho \partial^2_t \uu &= \nabla \cdot \stress + \mathbf{f} 
        && \text{Conservation de la quantité de mouvement} \label{eq:momentum}\\
    \strain &= \frac{1}{2}\left(\nabla\uu + \nabla^\intercal \uu\right)
        && \text{Tenseur des déformations infinitésimales} \label{eq:strain}\\
    \stress &= \mathbf{C} : \strain && \text{Équation constitutive} 
        \label{eq:stress}.
\end{align}

Les tenseurs symétriques de déformation $\strain$ et de contraintes $\stress$, 
sont liés par une loi constitutive à travers le tenseur des rigidités 
$\mathbf{C}$ d'ordre $4$. Dans le cas d'un matériau isotrope, on trouve la 
relation
\begin{align}
    \stress &= 2\mu \strain + \lambda \mathrm{tr}(\strain) \mathbf{I} 
    \label{eq:constitutive}\\
    \mu &= \frac{E}{2(1 + \nu)}\\
    \lambda &= \frac{E\nu}{(1+\nu)(1 - 2\nu)}
\end{align}

où 
\begin{itemize}
    \item $E$ [Pa], le module de Young, relie la déformation à la contrainte de 
    traction / compression.
    \item $\nu$ [-], le coefficient de Poisson, quantifie la déformation 
    d'un matériau dans la direction perpendiculaire à l'effort.
\end{itemize}

Bien que la physique du solide soit évidemment tri-dimensionnelle, il est 
souvent possible de réduire le problème à deux dimensions, voire une seule. Ici,
nous allons nous intéresser à des cas particuliers où l'on peut négliger
l'influence de la troisième dimension :
\begin{itemize}
    \item \textbf{Déformations planes} : la section d'un solide de grande 
    profondeur ne peut pas se déformer dans cette direction.
    \item \textbf{Tension planes} : un solide de faible épaisseur ne subit pas 
    de déformation dans la direction de l'épaisseur.
    \item \textbf{Axisymétrie} : un solide de révolution soumis à des forces
    axisymétriques se déforme indépendamment de la position angulaire.
\end{itemize}


\subsection{Déformations planes}
En considérant un déplacement uniquement planaire $\uu = u(x,y)\mathbf{e}_1 
+ v(x,y)\mathbf{e}_2$, dont on peut négliger la dépendance en $z$, on obtient 
les tenseurs de déformation, et de contrainte grâce à l'équation constitutive 
\eqref{eq:constitutive} :
\begin{gather}
    \strain =
    \begin{bmatrix}
        \epsilon_{xx} & \epsilon_{xy} & 0\\
        \epsilon_{xy} & \epsilon_{yy} & 0\\
        0 & 0 & 0
    \end{bmatrix} 
    =
    \begin{bmatrix}
        u_x & (u_x+v_y)/2 & 0\\
        (u_x+v_y)/2 & v_y & 0\\
        0 & 0 & 0
    \end{bmatrix} \label{eq:planeStrainStrain}
    \\[1em]
    \stress =
    \begin{bmatrix}
        \sigma_{xx} & \sigma_{xy} & 0\\
        \sigma_{xy} & \sigma_{yy} & 0\\
        0 & 0 & \sigma_{zz}
    \end{bmatrix}\\[1em]
    \begin{bmatrix}
        \sigma_{xx}\\
        \sigma_{yy}\\
        \sigma_{zz}\\
        \sigma_{xy}
    \end{bmatrix} = 
    \underbrace{
    \begin{bmatrix}
        \gamma & \lambda & \lambda & \\
        \lambda & \gamma & \lambda & \\
        \lambda & \lambda & \gamma & \\
         &  &  & \mu
    \end{bmatrix}}_{\mathbf{C}}
    \begin{bmatrix}
        \epsilon_{xx}\\
        \epsilon_{yy}\\[3pt]
        0\\
        2\epsilon_{xy}
    \end{bmatrix}\\[1em]
    \text{où} \quad
    \gamma \triangleq 2\mu + \lambda = 
    \frac{E(1 - \nu)}{(1 + \nu)(1 - 2\nu)}.
\end{gather}

\subsection{Tensions planes}
En considérant que la pièce est soumise à des contraintes dans le plan
$xy$ uniquement, on force $0=\sigma_{zz}$ :
\begin{equation}
    \begin{aligned}
        \sigma_{zz} &= \lambda \epsilon_{xx}+ \lambda\epsilon_{yy}
        + (2\mu+\lambda) \epsilon_{zz} = 0\\
        \iff \epsilon_{zz} &= -\frac{\nu}{1-\nu} (\epsilon_{xx} + \epsilon_{yy})
    \end{aligned}
\end{equation}

En considérant à nouveau un déplacement $\uu = u(x,y)\mathbf{e}_1 
+ v(x,y)\mathbf{e}_2$, on obtient

\begin{gather}
    \strain =
    \begin{bmatrix}
        \epsilon_{xx} & \epsilon_{xy} & 0\\
        \epsilon_{xy} & \epsilon_{yy} & 0\\
        0 & 0 & \epsilon_{zz}
    \end{bmatrix}
    =
    \begin{bmatrix}
        u_x & (u_x+v_y)/2 & 0\\
        (u_x+v_y)/2 & v_y & 0\\
        0 & 0 & -\frac{\nu}{1-\nu} (u_x + v_y)
    \end{bmatrix}\label{eq:planeStressStrain}
    \\[1em]
    \stress =
    \begin{bmatrix}
        \sigma_{xx} & \sigma_{xy} & 0\\
        \sigma_{xy} & \sigma_{yy} & 0\\
        0 & 0 & 0
    \end{bmatrix}\\[1em]
    \begin{bmatrix}
        \sigma_{xx}\\
        \sigma_{yy}\\
        \sigma_{xy}
    \end{bmatrix} = 
    \underbrace{
    \begin{bmatrix}
        \alpha & \beta & \\
        \beta & \alpha & \\
         &  & \mu
    \end{bmatrix}}_{\mathbf{C}}
    \begin{bmatrix}
        \epsilon_{xx}\\
        \epsilon_{yy}\\
        2\epsilon_{xy}
    \end{bmatrix} \qquad 
    \begin{aligned}
        \alpha &= \frac{E}{1 - \nu^2}\\
        \beta &= \frac{E\nu}{1 - \nu^2}\\
    \end{aligned}
\end{gather}

\subsection{Axisymétrie}
En considérant un déplacement axisymétrique 
$\uu = u(r,z)\mathbf{e}_r + v(r,z) \mathbf{e}_{\theta} + w(r,z)\mathbf{e}_z$,
et en utilisant le gradient en coordonnées cylindriques, on obtient les 
relations
\begin{gather}
    \strain =
    \begin{bmatrix}
        \epsilon_{rr} & \epsilon_{r\theta} & \epsilon_{rz}\\
        0 & \epsilon_{\theta \theta} & 0\\
        \epsilon_{rz} & 0 & \epsilon_{zz}
    \end{bmatrix}
    =
    \begin{bmatrix}
        u_r & (-v/r+v_r)/2 & (u_z + w_r)/2\\
        (-v/r+v_r)/2 & u/r & v_z / 2\\
        (u_z + w_r)/2 & v_z / 2 & w_z
    \end{bmatrix}\label{eq:axisymStrain}
    \\[1em]
    \stress =
    \begin{bmatrix}
        \sigma_{rr} & \sigma_{r\theta} & \sigma_{rz}\\
        \sigma_{r\theta} & \sigma_{\theta \theta} & \sigma_{\theta z}\\
        \sigma_{rz} & \sigma_{\theta z} & \sigma_{zz}
    \end{bmatrix}\\[1em]
    \begin{bmatrix}
        \sigma_{rr}\\
        \sigma_{\theta\theta}\\
        \sigma_{zz}\\
        \sigma_{rz}\\
        \sigma_{r\theta}\\
        \sigma_{\theta z}
    \end{bmatrix} = 
    \underbrace{
    \begin{bmatrix}
        \gamma & \lambda & \lambda &  &  & \\
        \lambda & \gamma & \lambda &  &  & \\
        \lambda & \lambda & \gamma &  &  & \\
         &  &  & \mu &  & \\
         &  &  &  & \mu & \\
         &  &  &  &  & \mu
    \end{bmatrix}}_{\mathbf{C}}
    \begin{bmatrix}
        \epsilon_{rr}\\
        \epsilon_{\theta\theta}\\
        \epsilon_{zz}\\
        2\epsilon_{rz}\\
        2\epsilon_{r\theta}\\
        2\epsilon_{\theta z}
    \end{bmatrix}.
\end{gather}

Cependant, la composante tangentielle $v$ du déplacement est 
supposée nulle dans le code, ce qui nous permet de réduire la matrice 
$\mathbf{C}$ à ses $4$ premières lignes et colonnes :
\begin{equation}
    \begin{bmatrix}
        \sigma_{rr}\\
        \sigma_{\theta\theta}\\
        \sigma_{zz}\\
        \sigma_{rz}
    \end{bmatrix} = 
    \underbrace{
    \begin{bmatrix}
        \gamma & \lambda & \lambda &  \\
        \lambda & \gamma & \lambda &  \\
        \lambda & \lambda & \gamma &  \\
         &  &  & \mu 
    \end{bmatrix}}_{\mathbf{C}}
    \begin{bmatrix}
        \epsilon_{rr}\\
        \epsilon_{\theta\theta}\\
        \epsilon_{zz}\\
        2\epsilon_{rz}
    \end{bmatrix}.
\end{equation}

\newpage
\section{Élements finis}

\subsection{Formulation faible}
Soit un corps $\Omega \subseteq \mathbb{R}^d$ avec une frontière $\Gamma =
\partial \Omega$ partitionnée en 
\begin{itemize}
    \setlength{\itemsep}{0pt}
    \item $\Gamma_D$ : déplacement imposé (Dirichlet),
    \item $\Gamma_N$ : contrainte imposée (Neumann),
    \item $\Gamma_R$ : relation contrainte-déplacement imposée (Robin).
\end{itemize}

Les contraintes peuvent être imposées sur le vecteur complet, qu’il s’agisse 
du déplacement ou de la contrainte, ou bien seulement sur une de ses 
composantes $x$, $y$ voire même selon un direction normale, tangente, 
ou arbitraire.
Notonts $\mathcal{M}(\xx)$ l'ensemble des déplacements $\uu(\xx)$ admissibles 
sur un point de $\Gamma_D$. Dans le cas d'un déplacement normal, on a par 
exemple
\begin{equation}
    \mathcal{M}(\xx) \triangleq \left\{ \uu(\xx) \in \mathbb{R}^d \quad 
        \text{t.q.} \quad \uu(\xx) \cdot \mathbf{n}(\xx) = 0\right\}.
\end{equation}

En suivant l'approche de Galerkin, nous cherchons une solution 
$\uu \in \mathcal{U}(\Omega)$, l'espace des solutions admissibles, 
telle que la formulation faible soit vérifiée 
$\forall \vv \in \mathcal{V}(\Omega)$, l’espace des fonctions test :
\begin{alignat}{4}
    \mathcal{U} &\triangleq \big\{ 
        \uu \in [H^1(\Omega)]^d \quad &&\text{t.q.} \quad 
        \uu(\xx) - \uu_D(\xx) &&\in \mathcal{M}(\xx) 
        \quad &&\forall \xx \in \Gamma_D
    \big\},\\
    \mathcal{V} &\triangleq \big\{ 
        \vv \in [H^1(\Omega)]^d \quad &&\text{t.q.} \quad 
        \vv(\xx) &&\in \mathcal{M}(\xx)
        \quad &&\forall \xx \in \Gamma_D
    \big\}.
\end{alignat}

En partant de l'équation de conservation de la quantité de mouvement, on
obtient :
\begin{equation}
    \begin{aligned}
        \int_\Omega \rho \partial^2_t \uu \cdot \vv \:\dd\Omega 
        &= 
        \int_\Omega (\nabla \cdot \stress) \cdot \vv \:\dd\Omega + 
        \int_\Omega \mathbf{f} \cdot \vv \:\dd\Omega\\
        &=
        \int_\Omega \left[\nabla \cdot (\stress \cdot \vv) 
        - \stress : \nabla \vv \right] \:\dd\Omega + 
        \int_\Omega \mathbf{f} \cdot \vv \:\dd\Omega\\
        &=
        -\int_\Omega \stress : \left[
            \strain(\vv) + \boldsymbol{\omega}(\vv))\right] \:\dd\Omega +
        \int_{\Gamma}  \mathbf{n} \cdot \stress \cdot \vv \:\dd\Gamma +
        \int_\Omega \mathbf{f} \cdot \vv \:\dd\Omega\\
        &= 
        -\int_\Omega \stress(\uu) : \strain(\vv) \:\dd\Omega +
        \int_\Omega \mathbf{f} \cdot \vv \:\dd\Omega +
        \int_{\Gamma_N \cup \Gamma_R}  \mathbf{t} \cdot \vv \:\dd\Gamma 
        \qquad \forall \vv \in \mathcal{V}(\Omega).
    \end{aligned}\label{eq:weak}
\end{equation}

\subsection{Formulation discrète}
On fait maintenant l'hypothèse que la solution $\uu$ peut être approchée par 
une approximation linéaire par morceaux $\uu^h$. Chaque morceau est un élément
(triangle / quadrilatère) d'un maillage $\mathcal{T}_h$ de $\Omega$, avec 
$n$ nœuds et $N$ éléments. On note $\phi_{j}(\xx)$ la fonction de
forme scalaire associée au n\oe ud $j$ et $\uu_j$ la valeur du 
déplacement au n\oe ud $j$. On peut donc écrire la solution approchée
\begin{equation}
    \uu^h(\xx) = \sum_{j=1}^{n} \phi_j(\xx) \uu_j.
\end{equation}

Étant donné qu'il y a $2n$ degrés de liberté, la base de l'espace discret
$\mathcal{U}^h$ contient $2n$ éléments :
\begin{equation}
    \begin{alignedat}{9}
    \mathcal{U}^h &=
    &&\Bigg\{
        \begin{bmatrix}
            \phi_1\\
            0
        \end{bmatrix}
        &&,
        \begin{bmatrix}
            0\\
            \phi_1
        \end{bmatrix}
        &&,
        \begin{bmatrix}
            \phi_2\\
            0
        \end{bmatrix}
        &&,
        \begin{bmatrix}
            0\\
            \phi_2
        \end{bmatrix}
        &&,
        \;\: \ldots\;\:
        &&,
        \begin{bmatrix}
            \phi_n\\
            0
        \end{bmatrix}
        &&,
        \begin{bmatrix}
            0\\
            \phi_n
        \end{bmatrix}
    &&\Bigg\}\\ 
    &=
    &&\Big\{
        \;\: \shapeu_{1} &&,
        \;\: \shapev_{1} &&,
        \;\: \shapeu_{2} &&,
        \;\: \shapev_{2} &&,
        \;\: \ldots \;\: &&,
        \;\: \shapeu_{n} &&,
        \;\: \shapev_{n}
    &&\Big\}\\ 
    &=
    &&\Big\{
        \;\: \shapef_{1}&&,
        \;\: \shapef_{2}&&,
        \;\: \shapef_{3}&&,
        \;\: \shapef_{4}&&,
        \;\: \ldots \;\: &&,
        \;\: \shapef_{2n-1}&&,
        \;\: \shapef_{2n}
    &&\Big\}
    \end{alignedat}
\end{equation}

On peut ainsi réécrire l'approximation de la solution $\uu^h$ comme
\begin{equation}
    \uu^h(\xx)
    = \sum_{j=1}^{n} \shapeu_{j}(\xx) u_j + \shapev_{j}(\xx) v_j    
    = \sum_{j=1}^{2n} \shapef_j(\xx) \: q_j
\end{equation}

En remplaçant $\uu$ par $\uu^h$ dans \eqref{eq:weak}, on obtient la
formulation discrète :
\begin{equation}
    \begin{aligned}
        \sum_{j=1}^{2n} \bigg[&\int_\Omega \rho \shapef_i \cdot \shapef_j \: \dd \Omega \: \bigg]
        \partial^2_t q_j\\
        + \sum_{j=1}^{2n} \bigg[&\int_\Omega \strain(\shapef_i) \cdot \mathbf{C} \cdot \strain(\shapef_j)\: \dd \Omega\bigg] q_j\\
        = &\int_\Omega \shapef_i \cdot \mathbf{f} \: \dd \Omega \: + \int_{\Gamma_N \cup \Gamma_R} \shapef_i \cdot
        \mathbf{t} \: \dd \Gamma
    \end{aligned},
    \label{eq:discrete}
\end{equation}

que l'on peut réécrire sous la forme matricielle
\begin{equation}
    M \ddot{q} + K q = b. \label{eq:matrix}
\end{equation}

Dans le cas plane, toutes les intégrales de volume présentées ci-dessus sont 
calculées dans un volume tri-dimensionel : 
$\int_{\Omega} = \int \int \int \dd x \dd y \dd z$.
Cependant, aucune intégrande ne dépend de la coordonnée $z$ et on peut donc
omettre une des intégrales des deux côtés de l'équation.

Dans le cas axisymétrique, les intégrales de volume sont également calculées 
dans un volume tri-dimensionel :
$\int_{\Omega} = \int \int \int r \: \dd r \dd \theta \dd z$.
À nouveau, on peut omettre l'intégrale sur $\theta$ car aucune intégrande ne
dépend de cette coordonnée.

\subsection{Matrices locales}

Les intégrales définies dans \eqref{eq:discrete} sont nulles presque partout, 
sauf lorsque $\shapef_i$ et $\shapef_j$ sont non nulles sur le même élément
$\Omega_e$. Cela se produit lorsque $i$ et $j$ sont des n\oe uds de l'élément
$e$. On assemble donc l'équation \eqref{eq:matrix} en ajoutant les contributions
de chaque élément $e$ du maillage $\mathcal{T}_h$.

Sur un élément $\Omega_e$ avec $m=3$ ou $m=4$ n\oe uds, on définit donc les 
matrices locales $M^e, K^e \in \mathbb{R}^{2m \times 2m}$ qui lient les degrés 
de liberté $q_i$ et $q_j$ de l'élément $e$, $1\leq i,j\leq m$. 
À chaque paire de n\oe ud $(i,j)$ correspond donc un bloc 
$2 \times 2$ $M_{ij}^e$ et $K_{ij}^e$ liant les deux degrés de liberté 
($u$, $v$) associés à ces n\oe uds.
\begin{align}
    M_{ij}^e &= 
    \begin{bmatrix}
        \langle \rho \shapeu_i \cdot \shapeu_j \rangle &
        \langle \rho \shapeu_i \cdot \shapev_j \rangle \\
        \langle \rho \shapev_i \cdot \shapeu_j \rangle &
        \langle \rho \shapev_i \cdot \shapev_j \rangle
    \end{bmatrix}
    = 
    \begin{bmatrix}
        \langle\rho \phi_i \phi_j\rangle & 0\\
        0 & \langle\rho \phi_i \phi_j\rangle
    \end{bmatrix} \label{eq:massMatrix}\\
    K_{ij}^e &= 
    \begin{bmatrix}
        \langle\strain(\shapeu_i) \cdot \mathbf{C} \cdot \strain(\shapeu_j)\rangle &
        \langle\strain(\shapeu_i) \cdot \mathbf{C} \cdot \strain(\shapev_j)\rangle\\
        \langle\strain(\shapev_i) \cdot \mathbf{C} \cdot \strain(\shapeu_j)\rangle &
        \langle\strain(\shapev_i) \cdot \mathbf{C} \cdot \strain(\shapev_j)\rangle
    \end{bmatrix}.\label{eq:stiffMatrix}
\end{align}


% L'expression de la matrice de raideur $K_{ij}^e$ peut être simplifiée grâce aux
% nombreuses entrées nulles présentes dans le tenseur de rigidité $\mathbf{C}$. 

Dans le code, on utilise la forme générale \eqref{eq:stiffMatrix} pour plus 
de clarté et de modularité. 
Cela nous permet de changer le modèle constitutif sans avoir à 
modifier le code source. En contre-partie, on est forcé de calculer les deux
produits matrices-vecteurs malgré les quelques entrées nulles.

Pour ce faire, il faut tout de même calculer les expressions de 
$\strain(\shapef_i)$.

\subsubsection{Cas plane}
Dans le cas des tensions planes et des déformations planes, le terme 
$\epsilon_{zz} \sigma_{zz}$ de la double contraction disparait, et on trouve
grâce à \eqref{eq:planeStrainStrain} que
\begin{equation}
    \strain(\shapeu_{i})_{\textrm{vec}} = 
    \begin{bmatrix}
        \epsilon_{xx} \\
        \epsilon_{yy} \\
        2\epsilon_{xy}
    \end{bmatrix} = 
    \begin{bmatrix}
        \phi_{i,x}\\
        0\\
        \phi_{i,y}
    \end{bmatrix} \qquad\qquad
    \strain(\shapev_{i})_{\textrm{vec}} = 
    \begin{bmatrix}
        \epsilon_{xx} \\
        \epsilon_{yy} \\
        2\epsilon_{xy}
    \end{bmatrix} = 
    \begin{bmatrix}
        0\\
        \phi_{i,y}\\
        \phi_{i,x}
    \end{bmatrix}.
\end{equation}

Ce qui nous permet d'écrire la matrice de raideur locale pour les déformations
planes :
\begin{equation}
    \begin{aligned}
        K_{ij}^e &= \int_{\Omega_e}
        \begin{bmatrix}
            \phi_{i,x} & 0 & \phi_{i,y}\\
            0 & \phi_{i,y} & \phi_{i,x}
        \end{bmatrix} 
        \begin{bmatrix}
            \gamma & \lambda & 0\\
            \lambda & \gamma & 0\\
            0 & 0 & \mu
        \end{bmatrix}
        \begin{bmatrix}
            \phi_{j,x} & 0\\
            0 & \phi_{j,y}\\
            \phi_{j,y} & \phi_{j,x}
        \end{bmatrix} \dd \Omega_e\\
        &=
        \left[
            \setlength{\arraycolsep}{12pt}
            \begin{array}{c|c}
                \langle\phi_{i,x} \:\gamma  \: \phi_{j,x}\rangle 
                + \langle\phi_{i,y} \: \mu \: \phi_{j,y}\rangle &
                \langle\phi_{i,x} \:\lambda \: \phi_{j,y}\rangle 
                + \langle\phi_{i,y} \: \mu \: \phi_{j,x}\rangle\\[1ex]
                \hline & \\[-1.5ex]
                \langle\phi_{i,y} \:\lambda \: \phi_{j,x}\rangle 
                + \langle\phi_{i,x} \: \mu \: \phi_{j,y}\rangle &
                \langle\phi_{i,y} \:\gamma  \: \phi_{j,y}\rangle 
                + \langle\phi_{i,x} \: \mu \: \phi_{j,x}\rangle
            \end{array}
        \right].
    \end{aligned} \label{eq:stiffMatrixPlane}
\end{equation}

Dans le cas des tensions planes, il suffit de remplacer $\gamma$ par $\alpha$ 
et $\lambda$ par $\beta$ dans \eqref{eq:stiffMatrixPlane}.

\subsubsection{Cas axisymétrique}
Dans le cas axisymétrique, on doit également prendre en compte la
composante $\epsilon_{\theta\theta}$ du tenseur de déformation comme mentionné
à l'équation \eqref{eq:axisymStrain} :
\begin{equation}
    \strain(\shapeu_{i})_{\textrm{vec}} = 
    \begin{bmatrix}
        \epsilon_{rr} \\
        \epsilon_{\theta\theta} \\
        \epsilon_{zz} \\
        2\epsilon_{rz}
    \end{bmatrix} =
    \begin{bmatrix}
        \phi_{i,r}\\
        \phi_i / r\\
        0\\
        \phi_{i,z}
    \end{bmatrix} \qquad\qquad
    \strain(\shapev_{i})_{\textrm{vec}} =
    \begin{bmatrix}
        \epsilon_{rr} \\
        \epsilon_{\theta\theta} \\
        \epsilon_{zz} \\
        2\epsilon_{rz}
    \end{bmatrix} =
    \begin{bmatrix}
        0\\
        0\\
        \phi_{i,z}\\
        \phi_{i,r}
    \end{bmatrix}.
\end{equation}

Il ne reste plus qu'à calculer la double contraction avec la matrice $C$ 
et à intégrer en coordonnées cylindriques :
\begin{equation}
    \begin{aligned}
        K_{ij}^e &= \int_{\Omega_e}
        \begin{bmatrix}
            \phi_{i,r} & \phi_i/r & 0 & \phi_{i,z}\\
            0 & 0 & \phi_{i,z} & \phi_{i,r}
        \end{bmatrix}
        \begin{bmatrix}
            \gamma & \lambda & \lambda & 0\\
            \lambda & \gamma & \lambda & 0\\
            \lambda & \lambda & \gamma & 0\\
            0 & 0 & 0 & \mu
        \end{bmatrix}
        \begin{bmatrix}
            \phi_{j,r} & 0\\
            \phi_{j}/r & 0\\
            0 & \phi_{j,z}\\
            \phi_{j,z} & \phi_{j,r}
        \end{bmatrix} \dd \Omega_e\\
        &=
        \left[
        \setlength{\arraycolsep}{6pt}
        \begin{array}{c|c}
            \langle\phi_{i,r} \:r\gamma  \: \phi_{j,r}\rangle + \langle\phi_{i,z} \: r\mu \: \phi_{j,z}\rangle &
            \langle\phi_{i,r} \:r\lambda \: \phi_{j,z}\rangle + \langle\phi_{i,z} \: r\mu \: \phi_{j,r}\rangle \\
            {} + \textcolor{Red}{\langle\phi_{i,r} \: \lambda \: \phi_j\rangle} + 
            \textcolor{Red}{\langle\phi_i \: \lambda \: \phi_{j,r}\rangle} & 
            {} + \textcolor{Red}{\langle\phi_i \: \lambda \: \phi_{j,z}\rangle}\\
            \textcolor{Red}{\langle\phi_i \: \gamma/r \: \phi_j\rangle} & \\[1ex]
            \hline & \\[-1.5ex]
            \langle\phi_{i,z} \:r\lambda \: \phi_{j,r}\rangle + \langle\phi_{i,r} \: r\mu \: \phi_{j,z}\rangle &
            \langle\phi_{i,z} \:r\gamma  \: \phi_{j,z}\rangle + \langle\phi_{i,r} \: r\mu \: \phi_{j,r}\rangle\\
            {} + \textcolor{Red}{\langle\phi_{i,z} \: \lambda \: \phi_j\rangle} & 
        \end{array}
        \right].
    \end{aligned}
\end{equation}

On remarque que les termes en noir sont similaires au cas plane alors que
des contributions supplémentaires apparaissent en rouge. Ces contributions sont
dues à la présence d'une déformation tangentielle $\epsilon_{\theta\theta}$ 
malgré l'absence de déplacement dans cette direction.

\subsubsection{Terme source local}
Le terme source est lui bien plus simple à assembler :
\begin{equation}
    b_{i}^e = 
    \begin{bmatrix}
        \langle\shapeu_i \cdot \mathbf{f}\rangle\\
        \langle\shapev_i \cdot \mathbf{f}\rangle
    \end{bmatrix} =
    \left\{
        \begin{array}{ll}
            \begin{bmatrix}
                \langle\phi_i \phantom{r} f_1\rangle\\
                \langle\phi_i \phantom{r} f_2\rangle
            \end{bmatrix} & \text{cas plane}
            \\[1em]
            \begin{bmatrix}
                \langle\phi_i r f_r\rangle\\
                \langle\phi_i r f_z\rangle
            \end{bmatrix} & \text{cas axisymétrique}
        \end{array}
    \right.
\end{equation}


\subsection{Conditions naturelles de Neumann et Robin}
Ces conditions sont appelées \textit{naturelles} car elles apparaissent
naturellement dans la formulation faible. Elles découlent de la contrainte 
imposée sur la frontière du solide :
\begin{alignat}{2}
    \mathbf{n} \cdot \sigma \vert_{\Gamma} &= \mathbf{t} 
        \qquad& \forall \xx &\in \Gamma_N\\
    \mathbf{n} \cdot \sigma \vert_{\Gamma} &= \mathbf{t} - 
        k (\mathbf{u} - \mathbf{u}_0) \qquad& \forall \xx &\in \Gamma_R.
\end{alignat}

\begin{itemize}
    \item Dans le cas de la condition de Neumann, on connait la contrainte au 
    bord.
    \item Dans le cas de la condition de Robin, on fait l'hypothèse que le bord 
    est attaché à un ressort de raideur $k$ autour d'un point de 
    référence $\xx+\mathbf{u}_0$.
\end{itemize}

On voit donc clairement que la condition de Neumann est un cas particulier de la
condition de Robin, lorsque $k=0$. Dans la suite, nous allons donc exclusivement 
nous intéresser à la condition de Robin sous la forme
\begin{equation}
    \mathbf{n} \cdot \stress \vert_{\Gamma} = \mathbf{p} - k \mathbf{u}.
\end{equation}

Le terme de surface de l'équation \eqref{eq:weak} s'écrit alors
\begin{equation}
    \int_{\Gamma} (\mathbf{n} \cdot \stress) \cdot \vv \:\dd\Gamma.
\end{equation}

Pour appliquer cette condition dans la direction normale ou tangentielle 
uniquement, il faut décomposer chaque vecteur dans ces deux directions. 
N'importe quelle autre paire de direction orthgonales est également possible.

\begin{equation}
    \begin{aligned}
        &\int_{\Gamma} \big[ 
            (\mathbf{n}\cdot \stress \cdot \mathbf{n}) \mathbf{n}+ 
            (\mathbf{n}\cdot \stress \cdot \mathbf{t}) \mathbf{t}
        \big] \cdot \big[
            (\vv \cdot \mathbf{n}) \mathbf{n} + 
            (\vv \cdot \mathbf{t}) \mathbf{t}
        \big] \:\dd\Gamma\\
        = &\int_{\Gamma} (p_n - k_n \uu \cdot \mathbf{n}) (\shapef \cdot \mathbf{n}) \: \dd \Gamma+
        \int_{\Gamma} (p_t - k_t \uu \cdot \mathbf{t}) (\shapef \cdot \mathbf{t}) \: \dd \Gamma
    \end{aligned}
\end{equation}

La modification à apporter au système $K q = b$ pour une condition de Robin 
dans la direction $\mathbf{n}$ s'obtient en utilisant successivement 
$\shapef = \shapeu_i = \phi_i \mathbf{e}_1$ et 
$\shapef = \shapev_i=\phi_i \mathbf{e}_2$ :
\begin{alignat}{2}
    K_{ij} &\leftarrow K_{ij} &&+ 
    \begin{bmatrix}
        \llangle k_n n_1 \phi_j \; \phi_i n_1 \rrangle & 
        \llangle k_n n_2 \phi_j \; \phi_i n_1 \rrangle
        \\
        \llangle k_n n_1 \phi_j \: \phi_i n_2 \rrangle & 
        \llangle k_n n_2 \phi_j \: \phi_i n_2 \rrangle
        \\
    \end{bmatrix},\\
    b_i &\leftarrow \;\; b_i &&+
    \begin{bmatrix}
        \llangle p_n \phi_i n_1 \rrangle\\
        \llangle p_n \phi_i n_2 \rrangle
    \end{bmatrix}
\end{alignat}

La modification pour la direction $\mathbf{t}$ s'obtient en 
remplaçant $\mathbf{n}$, $k_n$ et $p_n$ par $\mathbf{t}$, $k_t$ et $p_t$.
% Le cas d'une force dirigée selon un des axe se déduit trivialement de
% l'expression ci-dessus en prenant $\mathbf{n} = \mathbf{e}_1$ ou
% $\mathbf{n} = \mathbf{e}_2$.
L'extension au cas axisymétrique est immédiate : il suffit d'ajouter 
un facteur $r$ dans les intégrales de surface.


\subsection{Conditions essentielles de Dirichlet}
Ces conditions sont appelées \textit{essentielles} car elles doivent être
imposées à la formulation discrète en modifiant la matrice de raideur.

On souhaite pouvoir imposer à notre solide un déplacement dans une seule 
direction en lui laissant la liberté de se déplacer dans les autres directions
(l'autre direction en 2D). Dans le cas d'une pièce qui repose sur une 
surface plane, on veut par exemple interdire tout déplacement normal à cette 
surface mais lui laisser la liberté de glisser dessus.

Dans le cas où l'on souhaite imposer le déplacement complet, il suffit de le 
faire successivement dans chaque direction.

\subsubsection{Dans une direction alignée avec un axe}
\label{sec:dirichlet-aligned}

Comme précisé précédemment, sur la frontière de Dirichlet $\Gamma_D$, les 
fonctions tests $\shapef$ dans la direction du déplacement imposé sont nulles.
Prenons par exemple un déplacement $\bar u$ imposé selon la direction 
$\mathbf{e}_x$ au n\oe ud $i$. 
Il faut alors effectuer les opérations suivantes sur le degré de liberté $p$ 
correspondant
\begin{equation}
    \begin{alignedat}{3}
        K_{pq} &\gets 0 &&\forall q \neq p \quad && \iff\quad \shapeu_i = \mathbf{0}\\
        K_{pp} &\gets 1\\
        b_p &\gets \bar u\\
        K_{qp} &\gets 0 &&\forall q \neq p \\
        b_{q} &\gets b_q - K_{qp} \bar u \qquad &&\forall q \neq p\\
    \end{alignedat}
\end{equation}

Les deux dernières lignes ne sont pas obligatoires, mais elles permettent de
maintenir la symétrie de la matrice de raideur.

\subsubsection{Dans une direction abitraire}
Disons que l'on souhaite imposer au n\oe ud $i$ un déplacement $\bar u$ dans 
une direction $\mathbf{n}$ qui forme avec $\mathbf{t}$ une base orthonormée. 
Il faut alors modifier à la fois les inconnues du n\oe ud $i$ et les fonctions 
tests.

Soit $Q_{\textrm{loc}} \in \mathbb{R}^{2\times 2}$, la matrice qui transforme 
un déplacement normal / tangent vers un déplacement aligné avec $\mathbf{e}_1$ /
$\mathbf{e}_2$ :
\begin{equation}
    \begin{bmatrix}
        u_1\\
        u_2
    \end{bmatrix} = 
    \underbrace{
    \begin{bmatrix}
        n_1 & t_1\\
        n_2 & t_2
    \end{bmatrix}}_{Q_{\textrm{loc}}}
    \begin{bmatrix}
        u_n \\
        u_t
    \end{bmatrix} \quad\iff\quad
    \begin{bmatrix}
        u_n\\
        u_t
    \end{bmatrix} =
    \underbrace{
    \begin{bmatrix}
        n_1 & n_2\\
        t_1 & t_2
    \end{bmatrix}}_{Q_{\textrm{loc}}^{*}}
    \begin{bmatrix}
        u_1\\
        u_2
    \end{bmatrix}
\end{equation}

Étendons $Q_{\textrm{loc}}$ à une matrice $Q \in \mathbb{R}^{2n\times 2n}$
dont seul le bloc $2\times 2$ associé au n\oe ud $i$ diffère de l'identité :
\begin{equation}
    q = 
    \begin{bmatrix}
        u_1\\
        v_1\\
        \vdots\\
        u_i\\
        v_i\\
        \vdots
    \end{bmatrix} = 
    Q z \quad\iff\quad
    z =
    \begin{bmatrix}
        u_1\\
        v_1\\
        \vdots\\
        u_n\\
        v_t\\
        \vdots
    \end{bmatrix} = 
    Q^{*} q \quad\iff\quad
    Q = 
    \begin{bmatrix}
        1 & 0 & & & & \\
        0 & 1 & & & & \\
        & & \ddots & & & \\
        & & & n_1 & t_1 & \\
        & & & n_2 & t_2 & \\
        & & & & & \ddots\\
    \end{bmatrix}
\end{equation}

Il faut ensuite "tourner" la matrice de raideur $K$ et le vecteur de forçage $b$
pour les ramener dans le bon repère :
\begin{align}
    K q = b \quad\longrightarrow\quad Q^* K Q z &= Q^* b
\end{align}

On peut ensuite appliquer sur ce nouveau système la procédure présentée à la
section \ref{sec:dirichlet-aligned} pour imposer le déplacement $\bar u$ au
n\oe ud $i$ dans la direction $\mathbf{n}$. Enfin, pour passer de la solution
$z$ à la solution $q$, il faut appliquer la transformation inverse, de part
et d'autre pour maintenir la symétrie. C'est en pratique ce qui est implémenté
dans le code source.

Essayons néanmoins de décrire la procédure de manière plus formelle.
Définissons la matrice $C = I - \mathbf{e}_p \mathbf{e}_p^*$, une 
matrice de projection dont seule la première entrée du bloc $2\times 2$ 
associé au n\oe ud $i$ diffère de l'identité.

\begin{alignat}{3}
    && K x &= b\\
    &\longrightarrow &\quad Q^* K Q z &= Q^* b  &\quad&\text{Rotation}\\
    &\longrightarrow &\quad C Q^* K Q z &= C Q^* b &\quad&\text{Annulation}\\
    &\longrightarrow &\quad (C Q^* K Q C + \mathbf{e}_p \mathbf{e}_p^*) z 
        &= C Q^* b - C Q^*KQ \mathbf{e}_p \bar u + \mathbf{e}_p \bar u
        &\quad&\text{Symétrisation}\\
    &\longrightarrow &\quad Q(C Q^* K Q C + \mathbf{e}_p \mathbf{e}_p^*) Q^*x 
        &= Q C Q^* b - Q C Q^*KQ \mathbf{e}_p \bar u + Q \mathbf{e}_p \bar u
        &\quad&\text{Rotation'}
\end{alignat}

Le système ci-dessus peut être habilement réécrit comme
\begin{gather}
    (N + T K T) x = T b - T K \tilde n \bar u + \tilde n \bar u, \qquad 
    \text{avec }\\[1em]
    \begin{array}{c|c|c}
        T \triangleq Q C Q^* = &
        N \triangleq Q \mathbf{e}_p \mathbf{e}_p^* Q^* = &
        \tilde n \triangleq Q \mathbf{e}_p\\
        \begin{bmatrix}
            1 & 0 &        &        &        & \\
            0 & 1 &        &        &        & \\
              &   & \ddots &        &        & \\
              &   &        & t_1^2  & \:t_1\:t_2 & \\
              &   &        & \:t_1\:t_2 & t_2^2  & \\
              &   &        &        &        & \ddots
        \end{bmatrix}\quad&
        \begin{bmatrix}
            0 & 0 &        &        &        & \\
            0 & 0 &        &        &        & \\
              &   & \ddots &        &        & \\
              &   &        & n_1^2  & n_1n_2 & \\
              &   &        & n_1n_2 & n_2^2  & \\
              &   &        &        &        & \ddots
        \end{bmatrix}\quad&
        =\begin{bmatrix}
            0\\
            0\\
            \vdots\\
            n_1\\
            n_2\\
            \vdots
        \end{bmatrix}.
    \end{array}
\end{gather}

On constate fort heureusement qu'avec $n_1=1=-t_2$ et $n_2=0=t_1$, on retombe sur 
le cas particulier de la section \ref{sec:dirichlet-aligned}.

\subsection{Calcul des contraintes}
Lorsque le système est résolu, on obtient le déplacement $\mathbf{u}_j$ 
pour chaque n\oe ud $j$. On peut ensuite calculer aux points d'intégration
les déformations grâce à la relation \eqref{eq:strain} et enfin les contraintes 
grâce au modèle constitutif \eqref{eq:stress}. 

Étant donné qu'on utilise des éléments linéaires triangulaires ($\mathcal{P}_1$)
ou quadrangulaires ($\mathcal{Q}_1$), les déformations $\strain^h$ sont 
discontinues par éléments. Par conséquent les contraintes $\stress^h$ le sont
également.
Il est toutefois possible de reconstruire une approximation $\stress^c$ 
continue aux n\oe uds du maillage :
\begin{equation}
    \stress^c(\xx) = \sum_{j=1}^{n} \phi_j(\xx) \stress_j^h.
\end{equation}

Au moins deux méthodes sont possibles :
\begin{itemize}
    \item Moyenner autour de chaque n\oe ud les contraintes des éléments 
    adjacents, ou bien
    \item Minimiser l'intégrale du carré de la différence entre les contraintes
    discontinues $\stress^h$ et les contraintes continues $\stress^c$.
\end{itemize}

\subsubsection{Moyenne pondérée}
Cette première méthode a l'avantage d'être rapide, mais fournit de moins bons
résultats. En notant $\stress_j^c$ la valeur nodale du champ de contrainte
et $D_j$ l'ensemble des éléments adjacents au n\oe ud $j$, on a alors
\begin{equation}
    \stress^c_j = \frac{\displaystyle\int_{D_j} \phi_j \stress^h \dd \xx}{
        \displaystyle\int_{D_j} \phi_j \dd \xx}.
\end{equation}

\subsubsection{Moindre carrés}
Cette seconde méthode fournit de meilleurs résultats mais se révèle plus 
coûteuse étant donné qu'il faut résoudre un système linéaire.
On cherche ici à minimiser l'intégrale
\begin{equation}
    J(\sigma_1, \dots, \sigma_n) = 
        \int_{\Omega} \Big(
            \sigma^c\big(\sigma_1,\dots, \sigma_n\big) - \sigma^h
        \Big)^2 \dd \Omega
\end{equation}
pour chaque composante $\sigma^c$ et $\sigma^h$ du tenseur de contrainte.

On trouve les paramètres $\sigma_i$ du minimiseur $\sigma^c$ grâce à la 
condition de stationnarité
\begin{equation}
    \begin{alignedat}{2}
        &&\partial_{\sigma_i} J =
        2\int_{\Omega} \Big(
            \sigma^c\big(\sigma_1,\dots, \sigma_n\big) - \sigma^h
            \Big) \partial_{\sigma_i} \sigma^c \:\dd \Omega &= 0\\
        &\implies& \int_{\Omega} \Big(\sum_j \phi_j \sigma_j - \sigma^h\Big) \phi_i \:\dd \Omega &= 0\\
        &\implies& \sum_j \left[\int_{\Omega} \big(\phi_j \cdot \phi_i \big) \:\dd \Omega \right] \sigma_j &=
        \int_{\Omega} \big(\sigma^h \cdot \phi_i\big) \:\dd \Omega
    \end{alignedat}
\end{equation}
où l'on utilise la notation d'Einstein pour les indices. On retrouve
une matrice de masse similaire à celle définie dans \eqref{eq:massMatrix}, 
hormis que celle-ci ne porte que sur un champ scalaire. 
L'opération $M \sigma = b$ présentée ci-dessus n'est en fait qu'une
projection de la contrainte $\stress^h$ sur l'espace $\mathcal{U}^h$.


\subsection{Analyse modale}
En faisant l'hypothèse qu'on cherche une solution harmonique 
$\uu = \Phi e^{\imath \omega t}$, on peut réécrire l'équation 
\eqref{eq:matrix} sans forcage $b$ sous la forme d'un problème aux valeurs
propres généralisé :
\begin{equation}
    K \Phi = \omega^2 M \Phi,
\end{equation}
dont les solutions non-triviales $(\omega_i^2, \Phi_i)$ satisfont
\begin{align}
    \Phi_i K \Phi_j &= \omega_i^2 \delta_{ij} \quad\text{et}\\
    \Phi_i M \Phi_j &= \delta_{ij}.
\end{align}

C'est à dire que les vecteurs propres $\Phi_i$ sont orthonormés à travers la 
matrice de masse.

% Pour plus de clarté, les détails du raisonnement \eqref{eq:weak} sont listés
% ci-dessous :
% \begin{itemize}
%     \item Multiplication de l'équation de conservation de la quantité de 
%     mouvement par une fonction test $\vv$ et intégration sur le domaine $\Omega$.
%     \item Intégration par partie du terme $\nabla \cdot \stress \cdot \vv$.
%     \item Application du théorème de la divergence pour obtenir le terme de
%     surface, qui s'annule sur la frontière de Dirichlet $\Gamma_D$.
%     \item Décomposition du gradient de déplacement $\nabla \vv$ en partie 
%     symétrique $\strain$ et antisymétrique $\boldsymbol{\omega}$ qui disparait 
%     lors de la double contraction avec le tenseur de contrainte $\stress$.
% \end{itemize}


\section{Résolution numérique}

\subsection{Système linéaire}

\subsection{Mode propre}


\end{document}